\documentclass[12pt,a4paper]{ctexbook}
\usepackage{geometry}
\geometry{margin=2.2cm}
\usepackage{setspace}
\setstretch{1.35}
\usepackage{hyperref}
\title{全宋诗选·poet.song.0}
\author{古典文献整理}
\date{}
\begin{document}
\maketitle
\tableofcontents
\newpage
\section{日詩}
\textbf{宋太祖}\par\vspace{0.5em}
欲出未出光辣達，千山萬山如火發。\par
須臾走向天上來，逐却殘星趕却月。\par

\section{句}
\textbf{宋太祖}\par\vspace{0.5em}
未離海底千山黑，纔到天中萬國明。\par

\section{登戎州江樓閑望}
\textbf{幸夤遜}\par\vspace{0.5em}
滿目江山四望幽，白雲高卷嶂烟收。\par
日回禽影穿疏木，風遞猿聲入小樓。\par
遠岫似屏橫碧落，斷帆如葉截中流。\par

\section{雪}
\textbf{幸夤遜}\par\vspace{0.5em}
片片飛來靜又閑，樓頭江上復山前。\par
飄零盡日不歸去，帖破清光萬里天。\par

\section{雲}
\textbf{幸夤遜}\par\vspace{0.5em}
因登巨石知來處，勃勃元生綠蘚痕。\par
靜即等閑藏草木，動時頃刻徧乾坤。\par
橫天未必朋元惡，捧日還曾瑞至尊。\par
不獨朝朝在巫峽，楚王何事謾勞魂。\par

\section{句  其一}
\textbf{幸夤遜}\par\vspace{0.5em}
若教作鎮居中國，爭得泥金在泰山。\par

\section{句  其二}
\textbf{幸夤遜}\par\vspace{0.5em}
纔聞暖律先偷眼，既待和風始展眉。\par

\section{句  其三}
\textbf{幸夤遜}\par\vspace{0.5em}
嚼處春冰敲齒冷，嚥時雪液沃心寒。\par

\section{句  其四}
\textbf{幸夤遜}\par\vspace{0.5em}
蒙君知重惠瓊實，薄起金刀釘玉深。\par

\section{句  其五}
\textbf{幸夤遜}\par\vspace{0.5em}
深妝玉瓦平無壠，亂拂蘆花細有聲。\par

\section{句  其六}
\textbf{幸夤遜}\par\vspace{0.5em}
片逐銀蟾落醉觥。\par

\section{句  其七}
\textbf{幸夤遜}\par\vspace{0.5em}
巧剪銀花亂，輕飛玉葉狂。\par

\section{句  其八}
\textbf{幸夤遜}\par\vspace{0.5em}
寒艷芳姿色盡明。\par

\section{金陵覽古 秦淮}
\textbf{朱存}\par\vspace{0.5em}
一氣東南王斗牛，祖龍潜爲子孫憂。\par
金陵地脈何曾斷，不覺真人已姓劉。\par

\section{金陵覽古 石頭城}
\textbf{朱存}\par\vspace{0.5em}
五城樓雉各相望，山水英靈宅帝王。\par
此地定由天造險，古來長恃作金湯。\par

\section{金陵覽古 段石岡}
\textbf{朱存}\par\vspace{0.5em}
孫吳紀德舊刊碑，草沒蟠螭與伏龜。\par
惆悵岡頭三段石，至今猶似鼎分時。\par

\section{金陵覽古 北渠}
\textbf{朱存}\par\vspace{0.5em}
金殿分來玉砌流，黑龍湖徹鳳池頭。\par
後庭花落恩波斷，翻與南唐作御溝。\par

\section{金陵覽古 新亭}
\textbf{朱存}\par\vspace{0.5em}
滿目江山異洛陽，昔人何必重悲傷。\par
倘能戮力扶王室，當自新亭復故鄉。\par

\section{金陵覽古 天闕山}
\textbf{朱存}\par\vspace{0.5em}
牛頭天際碧凝嵐，王導無稽亦妄談。\par
若指遠山爲上闕，長安應合指終南。\par

\section{金陵覽古 東山}
\textbf{朱存}\par\vspace{0.5em}
鎮物高情濟世才，欲隨猿鶴老巖隈。\par
山花處處紅粧面，髣髴如初擁妓來。\par

\section{金陵覽古 烏衣巷}
\textbf{朱存}\par\vspace{0.5em}
人物風流往往非，空餘陋巷作烏衣。\par
舊時簾幕無從覓，只有年年社燕歸。\par

\section{金陵覽古 阿育王塔}
\textbf{朱存}\par\vspace{0.5em}
窣堵凝然鎮梵宮，舉頭層級在雲中。\par
金棺舍利藏何處，鐸繞危簷聲撼風。\par

\section{金陵覽古 烏衣巷}
\textbf{朱存}\par\vspace{0.5em}
閥閱淪亡梐枑移，年年舊燕亦雙歸。\par
茅簷葦箔無冠蓋，不見烏衣見白衣。\par

\section{金陵覽古 後湖}
\textbf{朱存}\par\vspace{0.5em}
雷轟疊鼓火翻旗，三翼翩翩試水師。\par
驚起黑龍眠不得，狂風猛雨下多時。\par

\section{金陵覽古 半陽湖}
\textbf{朱存}\par\vspace{0.5em}
江南龍節水爲鄉，水不純陰又半陽。\par
一片湖光共深淺，兩般泉脈異溫凉。\par

\section{金陵覽古 潮溝}
\textbf{朱存}\par\vspace{0.5em}
流水東西傍帝臺，六朝重爲兩朝開。\par
曾有鷁首知高下，莫問魚舟識去來。\par

\section{金陵覽古 直瀆}
\textbf{朱存}\par\vspace{0.5em}
晝役人功夜鬼功，陽開陰闔幾時終。\par
不聞擲土江中語，爭得盈流一水通。\par

\section{金陵覽古 運瀆}
\textbf{朱存}\par\vspace{0.5em}
舳艫銜尾日無虛，更鑿都城引漕渠。\par
何事餒來貪雀穀，不知留得幾年儲。\par

\section{金陵覽古 鳳凰臺}
\textbf{朱存}\par\vspace{0.5em}
竹影桐陰滿舊山，鳳凰多載不飛還。\par
登臺只有吹簫者，爭得和鳴墮世間。\par

\section{偈}
\textbf{釋德韶}\par\vspace{0.5em}
通玄峰頂，不是人門。\par
心外無法，滿目青山。\par

\section{頌}
\textbf{釋德韶}\par\vspace{0.5em}
暫下高峰已顯揚，般若圜通遍十方。\par
人天浩浩無差別，法界縱橫處處彰。\par

\section{偈}
\textbf{釋志端}\par\vspace{0.5em}
來年二月二，與汝暫相棄。\par
燒灰散長江，勿占檀那地。\par

\section{句}
\textbf{釋志端}\par\vspace{0.5em}
木馬走似烟，石人趁不及。\par

\section{春社日寄李學士}
\textbf{李濤}\par\vspace{0.5em}
社翁今日沒心情，爲乞治聾酒一瓶。\par
惱亂玉堂將欲徧，依稀巡到第三廳。\par

\section{題泥水關不動尊院}
\textbf{李濤}\par\vspace{0.5em}
走却坐禪客，移將不動尊。\par
世間顛倒事，八萬四千門。\par

\section{贈山翁}
\textbf{李濤}\par\vspace{0.5em}
松蘿深處住，閑野不生愁。\par
鳥語烟嵐靜，水聲門戶秋。\par
花間歸洞路，山下釣魚舟。\par
沐浴聖王化，自憐絲滿頭。\par

\section{題處士林亭}
\textbf{李濤}\par\vspace{0.5em}
帝里高人宅，蒼苔遶徑深。\par
卷簾山入戶，摘菓鳥移林。\par
石沼養龜水，月臺留客琴。\par
生涯一樽酒，名利不關心。\par

\section{雜詩四首  其一}
\textbf{李濤}\par\vspace{0.5em}
滔滔東流水，赴海無歸期。\par
亭亭右轉日，今茲復來茲。\par
人生寄天地，百年七十稀。\par
思慮復營營，恐爲達此嗤。\par

\section{雜詩四首  其二}
\textbf{李濤}\par\vspace{0.5em}
燕雀賀厦屋，蟲蟻慕腥羶。\par
野花引飛蝶，古木來新蟬。\par
秋天百蟲吟，春天百鳥喧。\par
翟公真闇事，喜慍見色言。\par

\section{雜詩四首  其三}
\textbf{李濤}\par\vspace{0.5em}
明主不棄士，我自志山林。\par
爵服豈無華，才疏力難任。\par
鳥向深山棲，魚由深淵沉。\par
吾亦愛吾廬，高歌復微吟。\par

\section{雜詩四首  其四}
\textbf{李濤}\par\vspace{0.5em}
深巖有老翁，龐眉鬚鬢雪。\par
夜半呼我名，授我微妙訣。\par
字畫古籀樣，體勢訛復缺。\par
雙眸忽炯炯，須臾竟披閱。\par
至今得其傳，心會口難說。\par

\section{句}
\textbf{李濤}\par\vspace{0.5em}
一言寤主寧復聽，三諫不從歸去來。\par
諫晉主不從作。\par

\section{俞公巖}
\textbf{陳摶}\par\vspace{0.5em}
萬事若在手，百年聊稱情。\par
他時南嶽去，記得此巖名。\par

\section{歸隠}
\textbf{陳摶}\par\vspace{0.5em}
十年蹤跡走紅塵，回首青山入夢頻。\par
紫陌縱榮爭及睡，朱門雖貴不如貧。\par
愁聞劍戟扶危主，悶見笙歌聒醉人。\par
携取舊書歸舊隠，野花啼鳥一般春。\par

\section{贈金勵睡詩  其一}
\textbf{陳摶}\par\vspace{0.5em}
常人無所重，惟睡乃爲重。\par
舉世皆爲息，魂離神不動。\par
覺來無所知，貪求心愈用。\par
堪笑塵中人，不知夢是夢。\par

\section{贈金勵睡詩  其二}
\textbf{陳摶}\par\vspace{0.5em}
至人本無夢，其夢本遊仙。\par
真人本無睡，睡則浮雲烟。\par
爐裏近爲藥，壺中別有天。\par
欲知睡夢裏，人間第一玄。\par

\section{題石水澗}
\textbf{陳摶}\par\vspace{0.5em}
銀河灑落翠光冷，一派回環湛晚暉。\par
幾恨却爲頑石礙，琉璃滑處玉花飛。\par

\section{冬日晚望}
\textbf{陳摶}\par\vspace{0.5em}
山鬼煖或呼，溪魚寒不跳。\par
晚景愈堪觀，危峰露殘照。\par

\section{西峰}
\textbf{陳摶}\par\vspace{0.5em}
爲愛西峰好，吟頭盡日昂。\par
巖花紅作陣，溪水綠成行。\par
幾夜礙新月，半山無夕陽。\par
寄言嘉遁客，此處是仙鄉。\par

\section{華山}
\textbf{陳摶}\par\vspace{0.5em}
半夜天香入巖谷，西風吹落嶺頭蓮。\par
空愛掌痕侵碧漢，無人曾歎巨靈仙。\par

\section{與毛女遊}
\textbf{陳摶}\par\vspace{0.5em}
藥苗不滿笥，又更上危顛。\par
回指歸去路，相將入翠烟。\par

\section{詠毛女}
\textbf{陳摶}\par\vspace{0.5em}
曾折松枝爲寶櫛，又編栗葉作羅襦。\par
有時問著秦宮事，笑撚仙花望太虛。\par

\section{對御歌}
\textbf{陳摶}\par\vspace{0.5em}
臣愛睡，臣愛睡。\par
不卧氈，不蓋被。\par
片石枕頭，蓑衣鋪地。\par
震雷掣電鬼神驚，臣當其時正酣睡。\par
閑思張良，悶想范蠡。\par
說甚孟德，休言劉備。\par
三四君子只是爭些閑氣，爭如臣向青山頂頭。\par
白雲堆裏，展開眉頭。\par
解放肚皮，但一覺睡。\par
管什玉兔東昇，紅輪西墜。\par

\section{辭上歸進詩}
\textbf{陳摶}\par\vspace{0.5em}
草澤吾皇詔，圖南摶姓陳。\par
三峰千載客，四海一閑人。\par
世態從來薄，詩情自得真。\par
乞全麞鹿性，何處不稱臣。\par

\section{石刻詩}
\textbf{陳摶}\par\vspace{0.5em}
我謂浮榮真是幻，醉來捨轡謁高公。\par
因聆玄論冥冥理，轉覺塵寰一夢中。\par

\section{贈張乖崖}
\textbf{陳摶}\par\vspace{0.5em}
自吳入蜀是尋常，歌舞筵中救火忙。\par
乞得金陵養閑散，也須多謝鬢邊瘡。\par

\section{詩一首}
\textbf{陳摶}\par\vspace{0.5em}
華陰高處是吾宮，出即凌空跨曉風。\par
臺殿不將金鎖閉，來時自有白雲封。\par

\section{喜英公大師挂錫太華}
\textbf{陳摶}\par\vspace{0.5em}
暗喜蓮峰作近隣，撥開雲霧見師頻。\par
有時問箇艱難字，便沐周旋說與人。\par
唐李監應留後跡，漢蔡邕想是前身。\par
堪嗟繼踵無徒弟，筆法收藏在渭濱。\par

\section{句  其一}
\textbf{陳摶}\par\vspace{0.5em}
蠶月桑葉青，鶑時柳花白。\par

\section{句  其二}
\textbf{陳摶}\par\vspace{0.5em}
萬頃白雲獨自有，一枝仙桂阿誰無。\par

\section{句}
\textbf{單氏}\par\vspace{0.5em}
帝王師不得，日月老應難。\par

\section{經方干舊居}
\textbf{釋乾康}\par\vspace{0.5em}
鏡湖中有月，處士後無人。\par
荻笋抽高節，鱸魚躍老鱗。\par

\section{賦殘雪}
\textbf{釋乾康}\par\vspace{0.5em}
六出奇花已住開，郡城相次見樓臺。\par
時人莫把和泥看，一片飛從天上來。\par

\section{詩一首}
\textbf{釋乾康}\par\vspace{0.5em}
隔岸紅塵忙似火，當軒青嶂冷如冰。\par
烹茶童子休相問，報道門前是衲僧。\par

\section{答僧問心要作偈}
\textbf{釋法騫}\par\vspace{0.5em}
昨日相逢敘起居，今朝相見事還如。\par
如何却覓呈心要，心要如何特地疏。\par

\section{答僧問佛法大意作偈}
\textbf{釋廣原}\par\vspace{0.5em}
剎剎現形儀，塵塵具覺知。\par
性源常鼓浪，不悟未曾移。\par

\section{占貯橘子}
\textbf{葉簡}\par\vspace{0.5em}
圓如珠，赤如丹，倘能擘破分喫了，爭不慚愧洞庭山。\par

\section{占巾子}
\textbf{葉簡}\par\vspace{0.5em}
近來好裹束，各自競尖新。\par
稱無二三兩，因何號一斤。\par

\section{占兩鷄子}
\textbf{葉簡}\par\vspace{0.5em}
此物不難知，一雄兼一雌。\par
請將打破看，方明混沌時。\par

\section{句}
\textbf{聶崇義}\par\vspace{0.5em}
勿笑有三耳，全勝畜二心。\par

\section{贈夢英大師}
\textbf{楊昭儉}\par\vspace{0.5em}
紀贈歌詩數百人，序師多藝各求新。\par
未言篆隸飛龍鳳，且說風騷感鬼神。\par
琴有古聲清耳目，鶴無凡態惹埃塵。\par
英公所學還如此，不錯承恩近紫宸。\par

\section{題家園}
\textbf{楊昭儉}\par\vspace{0.5em}
池蓮憔悴無顔色，園竹低垂減翠陰。\par
園竹池蓮莫惆悵，相看恰似主人心。\par

\section{自述真讚}
\textbf{釋通慧}\par\vspace{0.5em}
真兮寥廓，郢人圖雘。\par
嶽聳雲空，澄潭月躍。\par

\section{題玉堂壁}
\textbf{陶穀}\par\vspace{0.5em}
官職有來須與做，才能用處不憂無。\par
堪笑翰林陶學士，一生依樣畫葫蘆。\par

\section{石橋}
\textbf{陶穀}\par\vspace{0.5em}
重重翠幛聳雲端，玉殿金樓縹緲間。\par
聖境不容凡俗到，故將飛瀑隔塵寰。\par

\section{寄贈夢英大師}
\textbf{陶穀}\par\vspace{0.5em}
是箇碑文念得全，聰明靈性自天然。\par
離吳別楚三千里，入洛遊梁二十年。\par
負藝已聞喧世界，高眠長見卧雲烟。\par
相逢與我情何厚，問佛方知宿有緣。\par

\section{句  其一}
\textbf{陶穀}\par\vspace{0.5em}
尖簷帽子卑凡廝，短靿靴兒末厥兵。\par

\section{句  其二}
\textbf{陶穀}\par\vspace{0.5em}
此生頭已白，無路掃王門。\par

\section{句  其三}
\textbf{陶穀}\par\vspace{0.5em}
井蛙休恃重溟險，澤馬曾嘶九曲濱。\par

\section{句  其四}
\textbf{陶穀}\par\vspace{0.5em}
乞與金鐘病眼明。\par

\end{document}